# ResNet1D Baseline and Calibration Study on PTB-XL

## 1. Study objective

The objective of this phase is to establish a reproducible ResNet1D baseline for multi-label ECG classification on PTB-XL, evaluate alternative loss functions, and investigate whether post-hoc temperature scaling improves the reliability of predicted probabilities.

The study evaluates discrimination, threshold-dependent classification, calibration, and patient-level uncertainty. The ResNet1D experiments form the baseline for subsequent comparisons with alternative temporal architectures.

## 2. Dataset and preprocessing

- **Dataset:** PTB-XL
- **Recordings:** 21,799 in the local metadata
- **Sampling frequency:** 100 Hz
- **Input dimensions:** 12 ECG leads × 1,000 samples
- **Prediction task:** Multi-label classification across 44 labels
- **Training split:** 17,418 recordings
- **Validation split:** 2,183 recordings
- **Test split:** 2,198 recordings
- **Test patients:** 1,904 unique patients

The official patient-stratified folds are used: folds 1–8 for training, fold 9 for validation, and fold 10 for testing.

Preprocessing consists of third-order Butterworth filtering from 0.5 to 40 Hz, zero-phase filtering using `filtfilt`, and per-lead mean and standard-deviation normalization.

The prediction arrays have been verified against the test dataset targets in dataset order. All target entries match exactly.

## 3. Experimental configurations

Three loss configurations were evaluated:

1. **Standard binary cross-entropy (BCE):** Primary baseline.
2. **Weighted BCE:** Class-weighted loss intended to address label imbalance.
3. **Focal loss:** Gamma = 2, without class weights.

The baseline checkpoint is `checkpoints/best_model.pt`. The weighted-BCE and focal-loss checkpoints are `checkpoints/best_model_weighted_bce.pt` and `checkpoints/best_model_focal.pt`, respectively.

The primary model selection metric was validation macro-AUROC. Threshold-0.5 results and validation-selected threshold results are reported separately.

## 4. Test-set results

| Configuration | Best validation macro-AUROC | Test macro-AUROC | Test macro-AUPRC | Test macro-F1 at 0.5 |
|---|---:|---:|---:|---:|
| Standard BCE | 0.9260 | 0.9235 | 0.3041 | 0.2150 |
| Weighted BCE | 0.9270 | 0.9089 | 0.2932 | 0.2285 |
| Focal loss | 0.9223 | 0.9254 | Not recorded in the current consolidated results | 0.2097 |

The standard BCE model achieved a test macro-AUROC of 0.9235 and was retained as the primary baseline. Weighted BCE achieved a slightly higher validation macro-AUROC but lower test macro-AUROC. Focal loss achieved the highest recorded test macro-AUROC among these configurations, at 0.9254, but its threshold-0.5 macro-F1 was lower than the baseline.

These results illustrate that model ranking performance and threshold-dependent classification performance do not necessarily improve together. The existing experiments do not establish that one loss function is uniformly superior across all metrics.

### 4.1 Baseline performance by label-frequency tier

| Metric | Common | Medium | Rare |
|---|---:|---:|---:|
| AUROC | 0.9244 | 0.9086 | 0.9377 |
| AUPRC | 0.5795 | 0.2496 | 0.1017 |
| F1 at threshold 0.5 | 0.4984 | 0.1507 | 0.0148 |

Rare-label AUROC is comparatively high, but rare-label AUPRC and F1 remain low. This demonstrates why ranking metrics alone are insufficient to characterize multi-label classification performance, particularly for infrequent labels.

### 4.2 Threshold selection

For the baseline, validation-selected per-label thresholds produced:

- Macro-F1: 0.2853
- Micro-F1: 0.4425
- Common-tier F1: 0.5632
- Medium-tier F1: 0.2408
- Rare-tier F1: 0.0810

For weighted BCE, validation-selected thresholds produced macro-F1 of 0.2856 and micro-F1 of 0.5424.

Threshold-selected results must remain distinct from threshold-0.5 results. Thresholds should be selected using validation data only, and rare-label thresholds should be interpreted cautiously because they can be unstable when positive examples are scarce.

## 5. Calibration analysis

Global temperature scaling was fitted on validation predictions by minimizing binary cross-entropy. The fitted temperature was:

**T = 1.101489**

Validation BCE decreased from 0.063208 to 0.062705.

### 5.1 Test-set calibration results

| Metric | Raw probabilities | Temperature-scaled |
|---|---:|---:|
| Macro ECE | 0.010152 | 0.009712 |
| Common ECE | 0.025734 | 0.023787 |
| Medium ECE | 0.004930 | 0.004903 |
| Rare ECE | 0.000831 | 0.001385 |
| Macro Brier score | 0.017298 | 0.017225 |
| Common Brier score | 0.042746 | 0.042453 |
| Medium Brier score | 0.008843 | 0.008895 |
| Rare Brier score | 0.002001 | 0.002008 |

Temperature scaling preserves prediction ranking and leaves the observed AUROC, AUPRC, and threshold-0.5 F1 unchanged. It produces small reductions in aggregate ECE and Brier score, but the changes differ by frequency tier.

Common-label calibration improves on both ECE and Brier score. Rare-label ECE increases, and medium-label Brier score worsens slightly. Therefore, the results support a heterogeneous calibration effect rather than a universal improvement.

## 6. Patient-level uncertainty analysis

A paired patient-cluster bootstrap was performed using 1,000 replicates and random seed 42. The same patient resamples were used for raw and temperature-scaled predictions. The analysis includes 2,198 test recordings from 1,904 unique patients.

The saved macro-AUROC estimate is 0.923546, with a 95% percentile bootstrap interval of [0.910421, 0.931658]. Macro-F1 at threshold 0.5 is 0.214997, with an interval of [0.195094, 0.232371].

For calibration, the paired scaled-minus-raw differences and their existing percentile intervals are:

- **Macro ECE:** −0.000440; 95% CI [−0.000766, 0.000206].
- **Common ECE:** −0.001947; 95% CI [−0.002839, −0.000259].
- **Rare ECE:** +0.000554; 95% CI [+0.000083, +0.000863].
- **Macro Brier score:** −0.0000733; 95% CI [−0.0000983, −0.0000495].
- **Medium Brier score:** +0.0000519; 95% CI [+0.0000245, +0.0000782].

Negative differences favor temperature scaling for ECE and Brier score because lower values are better.

### 6.1 Statistical limitation

The observed full-test macro ECE lies below the lower endpoint of its bootstrap percentile interval. The same pattern appears for scaled macro ECE. A diagnostic reproduced this behavior, with no bootstrap replicate at or below the observed full-test estimate for either distribution.

This behavior warrants caution when interpreting the percentile intervals for ECE. It does not by itself prove an implementation error. The existing paired macro-ECE difference interval crosses zero, so the current analysis does not establish a clear aggregate macro-ECE improvement.

The intervals characterize test-sample uncertainty under patient-cluster resampling. They do not quantify variability across independent model-training seeds. The results should not be interpreted as evidence of clinical readiness.

## 7. Reproducibility artifacts

The following artifacts have been verified as present:

- `checkpoints/best_model.pt`
- `checkpoints/best_model_weighted_bce.pt`
- `checkpoints/best_model_focal.pt`
- `results/baseline_test_predictions.npz`
- `results/raw_calibration_test_predictions.npz`
- `results/temperature_scaled_test_predictions.npz`
- `results/temperature_scaling_results.json`
- `results/statistical_analysis_results.json`

The statistical results JSON parses successfully, and `src/statistical_analysis.py` compiles. Its interpretation note was corrected without recomputing or changing the numerical metrics.

## 8. Current conclusions and remaining limitations

The standard BCE ResNet1D provides a strong discrimination baseline on the selected PTB-XL test fold. Weighted BCE and focal loss offer useful ablations but show trade-offs between discrimination and threshold-dependent classification. Temperature scaling produces small, heterogeneous calibration changes.

The current record does not yet establish robustness across training seeds, clinical utility, or superiority over contemporary ECG architectures. Statistical significance should not be claimed beyond what the existing patient-level intervals support.

The proposed temporal-evidence reliability direction remains a research hypothesis, not a validated contribution. It will require a defined experimental protocol and a separate literature assessment before it can be claimed as novel.

## 9. Handoff to the next phase

The ResNet1D results are consolidated as the working baseline record. The next research phase is to implement ECG-Mamba as a separate model using the existing dataset splits, preprocessing, label vocabulary, and input/output contract.

The ECG-Mamba experiment should be evaluated against the frozen ResNet1D baseline using the same test protocol. Any additional temporal analysis should be motivated by a specific hypothesis rather than added solely to increase the number of experiments.